% Generated draft: requires mathematical and source review.
\documentclass[11pt]{article}
\usepackage[margin=25mm]{geometry}
\usepackage{fontspec}
\setmainfont{Latin Modern Roman}
\usepackage{amsmath,amssymb,mathtools,unicode-math}
\setmathfont{Latin Modern Math}
\usepackage{enumitem,needspace,xurl,hyperref}
\setlength{\parindent}{0pt}
\setlength{\parskip}{5pt}
\setlength{\emergencystretch}{3em}
\newcommand{\N}{\mathbb{N}}
\newcommand{\Z}{\mathbb{Z}}
\newcommand{\Q}{\mathbb{Q}}
\newcommand{\R}{\mathbb{R}}
\newcommand{\C}{\mathbb{C}}
\newcommand{\F}{\mathbb{F}}
\newcommand{\A}{\mathbb{A}}
\newcommand{\PP}{\mathbb{P}}
\newcommand{\E}{\mathbb{E}}
\newcommand{\eps}{\varepsilon}
\newcommand{\dd}{\,\mathrm{d}}
\begin{document}
\section*{UnsolvedMath ID: 800093}
\textbf{Source problem code:} catalogue entry
\subsection{Definitions and mathematical statement}
Let \(T^2=\R^2/\Z^2\), let \(S\) be a closed oriented surface of genus \(2\), and let \(\pi\colon M\to S\) be a fibration with fiber \(T^2\). Write \(\mathcal F\) for the foliation whose leaves are the fibers \(\pi^{-1}(s)\), for \(s\in S\).

The question is to identify hypotheses on the monodromy of \(\pi\) ensuring that every foliation \(\mathcal G\) sufficiently close to \(\mathcal F\) has a compact leaf close to some fiber of \(\pi\). In a locally trivial bundle interpretation, monodromy describes transport around loops in \(S\), up to the equivalence appropriate to the bundle category. The source does not specify that category or whether the intended hypotheses concern transport maps, their isotopy classes, or their induced action on the homology of \(T^2\).

The requested implication can be displayed as a schema. Let \(\mu_\pi\) denote the intended monodromy data, let \(\mathscr{Fol}(M)\) denote a chosen space of foliations of the same dimension as \(\mathcal F\), let \(\tau\) be a chosen topology on that space, and let \(C(L,\pi^{-1}(s))\) express the intended closeness of a leaf to a fiber. Seek a hypothesis \(H\) on \(\mu_\pi\) implying
\[
\begin{gathered}
\exists U\in\tau,\quad \mathcal F\in U,\\
\forall\mathcal G\in U\ \exists s\in S\ \exists L\text{ a leaf of }\mathcal G:\quad
L\text{ is compact}\ \text{ and }\ C(L,\pi^{-1}(s)).
\end{gathered}
\]
Here \(\mathscr{Fol}(M)\), \(\tau\), and \(C\) are placeholders for missing source specifications, not additional assumptions asserted by the source. In particular, the source does not specify differentiability, the topology of perturbation, the meaning of leaf closeness, or how the required leaf closeness depends on the size of the perturbation. Consequently this schema records the intended question but does not supply a uniquely determined precise formulation.
\subsection{Short English statement}
For a torus fibration over a closed oriented surface of genus \(2\), determine conditions on its monodromy that guarantee a compact leaf near a fiber after a sufficiently small perturbation of the fiber foliation. The source leaves the regularity, perturbation topology, leaf closeness, and intended monodromy data unspecified.
\subsection{Sources}
\begin{itemize}
\item Supplied local UnsolvedMath registry record, permanent ID \(800093\), source problem code \texttt{catalogue entry}, titled \emph{Foliations and Diffeomorphism Groups --- Question 34}. The record attributes the problem to Nariman and Yazdi, with proposed year \(2025\), and labels it \texttt{open}; that status is supplied metadata, not an independently verified mathematical assessment.
\item Nariman and Yazdi, \emph{Problem List on Foliations and Diffeomorphism Groups} (2025), Question 34. Supplied source URL: \url{https://arxiv.org/abs/2503.18274}. The supplied provenance records access on \texttt{2026-10-16} and extraction by \texttt{arxiv\_latex}; these are source-record statements, not a claim of independent access or verification.
\end{itemize}
\end{document}
